\section{Introduction}
\label{sec:intro}

\subsection{The question}

Two earlier papers \cite{Flat1644,Flat1671} studied a single fermion field with sixteen components on
two flat spacetimes of eight dimensions. One has four space-like and four time-like dimensions, and is
written $4+4$; the other has seven space-like dimensions and one time-like dimension, and is written
$7+1$. On both, the eight Dirac matrices can be chosen real, and the field supports a superfluid of
fermion pairs whose order parameter points along one of the four coordinates that are not the
coordinates of ordinary spacetime. \SIG{This paper treats the $4+4$ spacetime; a companion paper
\cite{CompanionDefl1671} treats $7+1$.}{This paper treats the $7+1$ spacetime; a companion paper
\cite{CompanionDefl1644} treats $4+4$.}

The eight coordinates are written $x0,x1,\dots,x7$. Throughout, $x1,x2,x3$ are the coordinates of
ordinary space and $x4$ is the ordinary time. The coordinate $x0$ is an additional space-like
coordinate, called the hidden coordinate. The remaining three, $x5,x6,x7$, are
\SIG{three additional \emph{time-like} coordinates}{three additional \emph{space-like} coordinates}.

In the flat theory nothing explains why the fields of ordinary physics do not depend on $x5$, $x6$
and $x7$. \SIG{On the $4+4$ spacetime this is more than a gap in the explanation. A field that
depends on an additional time does not evolve unitarily in the ordinary time, and the earlier paper
had to \emph{impose} that the fields do not depend on the additional times.}{On the $7+1$ spacetime
this was an assumption of the earlier paper.} The present paper examines a definite proposal for what
happens to these three coordinates.

\subsection{The hypothesis of deflation}

The hypothesis is that \emph{the coordinates $x5$, $x6$ and $x7$ deflate exponentially}: as the
ordinary time $x4$ advances, the proper extent of every interval along each of these coordinates
shrinks as an exponential of $x4$. Flat spacetime is replaced by a curved spacetime with a diagonal
metric, in which the three components of the metric along the deflating coordinates are
\begin{equation}
\begin{aligned}
g_{55}&=-\cff^{2}\,e^{-2\,\gff\,x4},\\
g_{66}&=-\css^{2}\,e^{-2\,\gss\,x4},\\
g_{77}&=-\cvv^{2}\,e^{-2\,\gvv\,x4}.
\end{aligned}
\label{eq:hyp}
\end{equation}
Here $\gff$, $\gss$ and $\gvv$ are three positive constants, the \emph{rates} of deflation; $\cff$,
$\css$ and $\cvv$ are three positive constants that set the scale at $x4=0$. The names of the rates
are written upright and without a subscript, to distinguish them from the components $g_{55}$,
$g_{66}$, $g_{77}$ of the metric, which carry subscripts. The other five components of the metric
keep their flat values. \SIG{The minus signs in Eq.~(\ref{eq:hyp}) are those of time-like components
in the convention of this paper, in which space-like components are positive.}{The minus signs in
Eq.~(\ref{eq:hyp}) are those of space-like components in the convention of this paper, in which the
time-like component is positive.}

The earlier papers examined three hypotheses about the superfluid, and we keep their labels.
\begin{itemize}
\item[H1.] The fermion field forms a superfluid state of spin 0 in which the fermions are paired by
the exchange of a new scalar boson, the \emph{Vorton-Scalar}, through a neutral current of the kind
known from the weak interaction.
\item[H2.] This superfluid state is the Higgs field.
\item[H3.] This superfluid state explains the observations of dark matter, or the Hubble tension, or
both.
\end{itemize}
The purpose of this paper is to repeat every calculation of the earlier paper under the hypothesis of
deflation, and to report what changes for H1, H2 and H3, favorable or not.

\subsection{Summary of results}

\emph{Geometry} (Sec.~\ref{sec:geometry}). The deflating spacetime is homogeneous and has 14
symmetries, or 15 or 20 if two or three of the rates are equal, instead of the 36 of flat spacetime.
It is not a solution of Einstein's equations without matter. With three equal rates it \emph{is} an
exact solution of the equations of Einstein-Gauss-Bonnet gravity without matter, for one definite
relation between the rate and the two constants of that theory.

\emph{The Dirac field} (Secs.~\ref{sec:dirac}, \ref{sec:modes}). A field that does not depend on
$x5,x6,x7$ obeys, after multiplication by the square root of the volume factor of the deflating
coordinates, exactly the flat equation. A field that does depend on them is affected directly.
%%<1644>
We solve its equation exactly. Of its sixteen solutions eight grow as the exponential of an
exponential of the time and are not admissible. The other eight are standing waves, which are
extinguished as the exponential of an exponential. \emph{The deflation removes the dependence of the
field on the additional times}, and with it the sector in which the $4+4$ theory is not unitary.
%%</1644>
%%<1671>
We solve its equation exactly. \emph{The deflation creates particles}: a mode that starts empty ends
with the occupation of a Fermi-Dirac distribution at a temperature equal to the rate divided by
$2\pi$, the Gibbons-Hawking temperature of the de Sitter space that the deflating coordinates and the
time form when the rates are equal.
%%</1671>

\emph{The superfluid} (Secs.~\ref{sec:clock} to \ref{sec:order}). For fields that do not depend on
the extra coordinates the model of the earlier paper holds unchanged, with one difference: its
coupling grows with the time, because a coupling of the four-dimensional description is a coupling
of the eight-dimensional theory divided by a power of the deflating volume. The masses do not change.
Every threshold of the earlier paper is therefore passed at a definite time. The gap of the paired
state grows much faster than the coupling. Two fermions become bound at a scattering resonance. The
paired state can follow the deflation only once its gap exceeds a few times the rate divided by the
pairing strength; before that it lags behind and breaks pairs. Matter of a given density loses its
fermion mass at a definite time. A vortex line carries four fermion zero modes.

\emph{Cosmology} (Sec.~\ref{sec:cosmo}). If gravity is the Einstein theory in eight dimensions,
Newton's constant grows at the total rate of deflation; lunar laser ranging then bounds that rate to
$3\times10^{-3}$ of the Hubble rate, and the Hubble constant inferred from the microwave background
falls. The deflation does not relieve the Hubble tension. Whatever keeps the rates constant must
carry a tension and must take part in the perturbations of ordinary space; the spectra of the
microwave background at all but the largest angular scales do not depend on how it does so, and
the growth of structure does.

\emph{The Higgs hypothesis} (Sec.~\ref{sec:higgs}). \SIG{The symmetry under which the pair condensate
has the pattern of the Higgs vacuum was a hidden symmetry already on the flat $4+4$ spacetime, and
remains one.}{The rotations of the four extra space-like coordinates, under which the pair condensate
has the pattern of the Higgs vacuum, are no longer symmetries of the spacetime; the pattern survives
only as a hidden symmetry.} If the condensate were the Higgs field, the electroweak scale would follow
the growing coupling, and atomic clocks would limit the rate of deflation to about $10^{-6}$ of the
Hubble rate. The hypothesis H2 and a deflation of any consequence today exclude each other.

\subsection{Method, and the limits of what is claimed}

Every statement of this paper about geometry, matrices and operators is an identity that a program
has verified in exact rational arithmetic, and that an independent program has verified again by
computer algebra for arbitrary values of the rates. Every numerical statement is the output of a
program, with a stated accuracy. The programs, their output and the documents that derive each result
step by step are public \cite{DeflationRepo}; the program performs {{CHECKS}} checks, all of which
pass. The metric of Eq.~(\ref{eq:hyp}) is a hypothesis. The paper derives its consequences; it does
not claim that the hypothesis is true. Section~\ref{sec:status} lists what has been established and
what has not.

Units are natural, $\hbar=c=k_B=1$. Symbols that carry the tag $\tg$ refer to the deflating
\SIG{$4+4$}{$7+1$} spacetime. A general coordinate among the three deflating ones is written $xt$ with
$t\in\{5,6,7\}$ and its rate $\rt$; a general coordinate among $x0,x1,x2,x3$ is written $xj$ with $j\in\{0,1,2,3\}$. The letter $a$ is kept for the four extra axes, $a\in\{0,5,6,7\}$ (Sec.~\ref{sec:gap}).
Appendix~\ref{app:symbols} lists the symbols.

